\subsection*{Resultados}

\begin{enumerate}[label=\alph*)]

\item \textbf{Construcción de la matriz de diseño $X$ y el vector de respuesta $Y$}

A partir de los datos experimentales, se construye la matriz de diseño incluyendo una columna de unos para el intercepto:

\[
X =
\begin{bmatrix}
1 & 1 & 2 \\
1 & 2 & 1 \\
1 & 2 & 3 \\
1 & 3 & 2
\end{bmatrix}
\]

y el vector de respuesta:

\[
Y =
\begin{bmatrix}
5 \\
4 \\
7 \\
6
\end{bmatrix}
\]

\item \textbf{Cálculo de la matriz transpuesta $X^T$}

La matriz transpuesta de $X$ es:

\[
X^T =
\begin{bmatrix}
1 & 1 & 1 & 1 \\
1 & 2 & 2 & 3 \\
2 & 1 & 3 & 2
\end{bmatrix}
\]

\item \textbf{Cálculo de $(X^T X)$ y $(X^T Y)$}

El producto matricial $X^T X$ es:

\[
X^T X =
\begin{bmatrix}
4 & 8 & 8 \\
8 & 18 & 16 \\
8 & 16 & 18
\end{bmatrix}
\]

El producto $X^T Y$ es:

\[
X^T Y =
\begin{bmatrix}
22 \\
47 \\
45
\end{bmatrix}
\]

\item \textbf{Planteamiento del sistema de ecuaciones}

El sistema matricial para estimar los parámetros del modelo es:

\[
(X^T X)\hat{\beta} = X^T Y
\]

Sustituyendo los valores obtenidos:

\[
\begin{bmatrix}
4 & 8 & 8 \\
8 & 18 & 16 \\
8 & 16 & 18
\end{bmatrix}
\begin{bmatrix}
\beta_0 \\
\beta_1 \\
\beta_2
\end{bmatrix}
=
\begin{bmatrix}
22 \\
47 \\
45
\end{bmatrix}
\]

Este sistema permite determinar los coeficientes del modelo de regresión lineal múltiple.

\end{enumerate}